\section{Exact rational evaluation and finite certificates}
\label{ado:sec:computation}
The density zeros need not be computed in order to obtain numerical
certificates. The unsigned norm bound $\norm\kappa_1\le2^{d-1}$
already gives an explicit geometric remainder for a centered
resolvent expansion. The proof works with real final order; the delivered
exact implementation accepts positive integer indices so that all
required moments are rational.

\subsection{A centered-moment expansion}
Define
\begin{equation}\label{ado:eq:centeredmoments}
 B_k=\int_0^1(2u-1)^k\kappa(u)\dd u
 =\sum_{j=0}^k(-1)^{k-j}\binom{k}{j}2^j c_{\svec}(j+1).
\end{equation}
Then $|B_k|\le M$, where $M=2^{d-1}$. For $\Ren z<1$, put
\[
 w=\frac{z}{2-z},\qquad |w|<1.
\]
The inequality is equivalent to $|z|<|2-z|$. The identity
$1-zu=(1-z/2)(1-w(2u-1))$ gives
\begin{theorem}[Geometric exact-evaluation certificate]
\label{ado:thm:evaluator}
For every integer $N\ge1$,
\begin{align}
 V_N(z)&=\frac{z}{1-z/2}\sum_{k=0}^{N-1}B_kw^k,
 \label{ado:eq:centeredvalue}\\
 |L(z)-V_N(z)|
 &\le\left|\frac{z}{1-z/2}\right|
          \frac{M|w|^N}{1-|w|}.
 \label{ado:eq:centerederror}
\end{align}
At $z=\ii$ and even $N=2m$, a purely rational bound is
\begin{equation}\label{ado:eq:gaussianerror}
 |L(\ii)-V_{2m}(\ii)|<\frac{2^d}{5^m}.
\end{equation}
For $|z|\le1/2$, the uniform bound is
\begin{equation}\label{ado:eq:halferror}
 |L(z)-V_N(z)|\le\frac{2^{d-1}}{3^N}.
\end{equation}
\end{theorem}
\begin{proof}
The centered geometric series converges uniformly for $u\in[0,1]$,
and its tail is bounded by $|w|^N/(1-|w|)$. Integrating against
$|\kappa|$ proves \eqref{ado:eq:centerederror}. At $z=\ii$,
$|w|=1/\sqrt5$ and $|z/(1-z/2)|=2/\sqrt5$. The multiplicative
constant is therefore $2M/(\sqrt5-1)<2M=2^d$ after extracting
$5^{-m}$. On $|z|\le1/2$, one has $|w|\le1/3$ and
$|z/(1-z/2)|\le2/3$, proving \eqref{ado:eq:halferror}.
\end{proof}
For integer indices, \eqref{ado:eq:centeredmoments} and
\eqref{ado:eq:centeredvalue} use only rational arithmetic at rational
complex $z$. An enclosing complex disk gives enclosing intervals for
each coordinate by adding and subtracting its radius. The evaluator
never treats a small residual as a proof of an identity.

\subsection{Implementation choices and replay}
The file \texttt{code/exact\_evaluator.py} first obtains the moments
by dynamic strict prefix sums. It computes $B_k$ with a common integer
denominator and a repeated-difference transform, rather than a
floating-point alternating binomial sum. Horner evaluation then gives
the exact rational center. The generic branch can upper-bound square
roots by dyadic rationals; the Gaussian and half-disk cases use the
closed rational bounds above. Near $\Ren z=1$, the fixed norm-bound
precision in the generic branch can be insufficient, in which case the
implementation raises an error rather than returns an uncertified value.
It is not advertised as a complete optimal evaluator on the entire slit
plane.

From the package root, the principal replay is
\begin{verbatim}
python code/verify.py
\end{verbatim}
It reads frozen rational certificates and independently recomputes them.
The optional command
\begin{verbatim}
python code/verify.py --regenerate
\end{verbatim}
uses floating point to \emph{propose} angular brackets, then requires
rational arithmetic to certify both endpoint signs before saving them.
The proposal stage uses a strict-disk series and does not decide any
accepted sign. This discovery/replay separation is deliberate.

\subsection{Gaussian values}
The delivered Gaussian certificate file contains eight mixed-depth
examples, all evaluated with $N=240$ centered terms. Each complex-disk
radius is smaller than $10^{-81}$. The decimal entries below are
readable midpoint diagnostics; the certificates are the rational
endpoints in \texttt{data/gaussian\_enclosures.json}.
\begin{center}\small
\begin{tabular}{@{}lrr@{}}
\toprule
Index & $\Ren L_{\svec}(\ii)$, approximately
      & $\Imn L_{\svec}(\ii)$, approximately\\
\midrule
$(2,1)$ & $-0.174715661158$ & $-0.113648917600$\\
$(1,1,1)$ & $0.099953993665$ & $-0.033577147847$\\
$(2,1,2)$ & $0.032383901801$ & $-0.026348629072$\\
$(1,2,1,1)$ & $0.001158016810$ & $0.008343470492$\\
$(2,1,1,2)$ & $0.001825088603$ & $0.006510815484$\\
$(1,1,1,1,1)$ & $-0.003396506582$ & $-0.001886779708$\\
$(2,1,2,1,1)$ & $-0.000326599679$ & $-0.000010475684$\\
$(1,1,1,1,1,1)$ & $0.000443168800$ & $-0.000335617002$\\
\bottomrule
\end{tabular}
\end{center}
In addition, the eight all-ones Gaussian values at depths one through
eight are checked against \emph{independent exact} intervals for
\[
 L_{\one_d}(\ii)=\frac{(-\tfrac12\log2+\tfrac{\ii\pi}{4})^d}{d!}.
\]
The logarithm and $\pi$ intervals are derived from different rational
series, as detailed in Appendix~\ref{ado:app:controls}. These eight
normalization controls are distinct from the eight stored mixed-depth
examples above.

\subsection{Certified angular brackets}
On the semicircle of radius $1/2$, set
\begin{equation}\label{ado:eq:rationalcircle}
 z(t)=\frac{1-t^2}{2(1+t^2)}+\ii\frac{t}{1+t^2},
 \qquad \theta=2\arctan t,\quad t>0.
\end{equation}
Rational $t$ gives rational real and imaginary coordinates. Every saved
root bracket has $t$-width at most $10^{-12}$; the corresponding
angle-width is less than $2\cdot10^{-12}$. With $N=104$, the
error bound \eqref{ado:eq:halferror} proves the two endpoint signs by
strict rational comparisons. For illustration:
\begin{center}\small
\begin{tabular}{@{}lcrr@{}}
\toprule
Index & $k$ & $10^{12}t_{\rm left}$ & $10^{12}t_{\rm right}$\\
\midrule
$(2,1,2)$ & 1 & 472099547153 & 472099547154\\
$(2,1,2)$ & 2 & 1487963538987 & 1487963538988\\
$(1,2,1,1)$ & 1 & 314972569528 & 314972569529\\
$(1,2,1,1)$ & 2 & 806524527583 & 806524527584\\
$(1,2,1,1)$ & 3 & 2019771345242 & 2019771345243\\
$(1,1,1,1,1,1)$ & 3 & 774596669241 & 774596669242\\
\bottomrule
\end{tabular}
\end{center}
The file \texttt{data/angular\_brackets.json} contains 24 brackets in
eight complete index families of depths two through six, together with
all exact coordinate intervals. The existence of a zero in each bracket
follows from the opposite signs. The theorem, rather than a scan,
proves uniqueness in every bracket and excludes all additional angular
zeros. For the displayed middle branch at depth six, the all-ones
formula independently gives
$\theta_3(1/2)=\arccos(1/4)$ and $t=\sqrt{3/5}$.

\subsection{What the exact suite checks}
The fresh replay reports 5,713 finite check groups:
\begin{center}\small
\begin{tabularx}{\linewidth}{@{}Xr@{}}
\toprule
Check family & Count\\
\midrule
Strict coefficients against direct finite nested enumeration & 1080\\
Initial zero moments and leading coefficients & 120\\
Centered moments: independent binomial formula and norm bound & 2880\\
Polynomial compensation identities for arbitrary test roots & 1200\\
All-ones moments against unsigned Stirling recurrence & 240\\
Logistic-kernel derivatives through depth twelve & 11\\
Depth-three/four symbolic kernel formulas & 2\\
Exact generalized-kernel counterexample checks & 2\\
Rejected negative controls & 2\\
Rational atomic radial-argument derivative signs & 80\\
Independent exact all-ones Gaussian controls & 8\\
Replayed stored Gaussian enclosures & 8\\
Replayed angular endpoint signs & 48\\
Replayed root brackets and complete families & $24+8$\\
\midrule
Total & 5713\\
\bottomrule
\end{tabularx}
\end{center}
The polynomial compensation tests deliberately use arbitrary rational
roots: they test the cancellation identity, \emph{not} positivity for
those roots. Positivity requires the actual density zeros and is proved
analytically. Similarly, finite atomic derivative checks do not replace
the interlacing-and-approximation proof of
Lemma~\ref{ado:lem:radialPick}.
